\chapter{Ouptut Files Format}
This chapter describes the format of the result files generated by the simulation process either in the \file{output} of the \file{misc} folders. The files in the \file{output} folder are text file whereas the \file{misc} folder contains binary files.

%%%%%%%%%%%%%%%%%%%%%%%%%
\section{output/ResultBoundaryFluxes.txt}
\index{Boundary}
%
\subsection{Purpose}
This file stores the instantaneous discharge crossing each boundary every \param{dt\_map}. In case of the boundary is composed of several interfaces, the stored discharge is the sum of the discharge crossing each interface. As the unit vector normal to the boundary interfaces is oriented from the exterior to the interior, a positive discharge enters the domain.
%
\subsection{Format description}
\texttt{<Tab>}: tabulation character. \\
    \begin{tabular}{p{6cm} p{9cm}}
    \hline
    Headerline & time \texttt{<Tab>} BC$_{x1}$\texttt{<Tab>}\ldots\texttt{<Tab>} BC$_{xN}$ where x1, \ldots, xN represent the nodestring numbers as defined in the \file{2dm file} \\
    \hline 
    $t$ \texttt{<Tab>} $Q_1$ \texttt{<Tab>} \ldots \texttt{<Tab>} $Q_N$ & $t$ is the stored time, $Q_i$ is the discharge (in $\mathrm{m^3.s^{-1}}$) through the boundary xN \\
    \hline
    \ldots & one line per stored time \\ \hline
    \end{tabular}
%
\subsection{Example}
\begin{verbatim}
time	BC_1	BC_2	BC_3
2.0000000000e+1	0.0000000000e+0	4.3274586317e+1	-1.4193271476e+1
3.0000000000e+1	0.0000000000e+0	4.3149535653e+1	-1.4227591312e+1
4.0000000000e+1	0.0000000000e+0	4.3074013841e+1	-1.4246148173e+1
5.0000000000e+1	0.0000000000e+0	4.3030038602e+1	-1.4258119344e+1
6.0000000000e+1	0.0000000000e+0	4.3000513361e+1	-1.4266604580e+1
\end{verbatim}

%%%%%%%%%%%%%%%%%%%%%%%%%
\section{output/ResultHydro\_pDDDDdHHhMMmSSsXXX.txt}

\subsection{Purpose}
This file stores the hydrodynamic results every \param{dt\_map}. To avoid any computational timestep reduction to store the results precisely every \param{dt\_map}, the results are stored when \pfg{explain algorithm}. The stored variables depend upon the chosen hyperbolic model (defined by the parameter \param{model\_hyperb}) and the \param{storeflag} parameter:
\begin{center}
\begin{tabular}{c|c|c}
\multirow{2}{*}{\param{model\_hyperb}} & \multicolumn{2}{c}{\param{storeflag} value} \\
& default & all \\
\hline
ddp  &  &  \\
dip  &  &  \\
sp   &  &  \\
swes & $h, q, r$ & 
\end{tabular}
\pfg{complete based on issue \#265}
\end{center}

\subsection{Format description}
The precise simulation time for storing the variables is indicated in the file name using the ISO 8601 format. The file name is \file{ResultHydro\_pDDDDdHHhMMmSSsXXX.txt} where \emph{pDDDDdHHhMMmSSsXXX} is the elapsed time from the beginning of the simulation:
\begin{itemize}
    \item DDDD is the number of days (always on 4 digits)
    \item HH is the number of hours (always on 2 digits)
    \item MM is the number of minutes (always on 2 digits)
    \item SS is the number of seconds (always on 2 digits)
    \item XXX is the number of milliseconds (always on 3 digits)
\end{itemize}

% \subsection{Example}
% \begin{verbatim}
    
% \end{verbatim}

% \begin{comment}
% %%%%%%%%%%%%%%%%%%%%%%%%%
% \section{output/ResultTransportReaction\_.txt}
% \index{Boussinesq coefficient}
% %
% \subsection{Purpose}
% %
% \subsection{Format description}
% %
% \subsection{Example}
% \begin{verbatim}
    
% \end{verbatim}
% \end{comment}

% %%%%%%%%%%%%%%%%%%%%%%%%%
% \section{output/ProbesTimeSeries.txt}
% \index{Probes}
% %
% \subsection{Purpose}
% %
% \subsection{Format description}
% %
% \subsection{Example}
% \begin{verbatim}
    
% \end{verbatim}

% %%%%%%%%%%%%%%%%%%%%%%%%%
% \section{output/ResultTime\_.txt}
% \index{Probes}
% %
% \subsection{Purpose}
% %
% \subsection{Format description}
% %
% \subsection{Example}
% \begin{verbatim}
    
% \end{verbatim}


% %%%%%%%%%%%%%%%%%%%%%%%%%
% \section{misc/ResultHydro\_}
% \index{Probes}
% %
% \subsection{Purpose}
% %
% \subsection{Format description}
% %
% \subsection{Example}
% \begin{verbatim}
    
% \end{verbatim}

% %%%%%%%%%%%%%%%%%%%%%%%%%
% \section{misc/ResultHydro.res}
% \index{Probes}
% %
% \subsection{Purpose}
% %
% \subsection{Format description}
% %
% \subsection{Example}
% \begin{verbatim}
    
% \end{verbatim}

% %%%%%%%%%%%%%%%%%%%%%%%%%
% \section{misc/ResultTransportReaction.res}
% \index{Probes}
% %
% \subsection{Purpose}
% %
% \subsection{Format description}
% %
% \subsection{Example}
% \begin{verbatim}
    
% \end{verbatim}
